\documentclass[12pt,oneside]{article}
\usepackage{listings}
\usepackage[ruled]{algorithm2e}
\usepackage[czech]{babel}
\usepackage[shortlabels]{enumitem}


\usepackage{titling}
\usepackage[a4paper,margin=2.5cm,head=0.5cm,foot=0.5cm]{geometry}
\usepackage{amsmath}
\usepackage{amsfonts}
\usepackage{amssymb}

\title{Domácí úkol k cvičení číslo 3}
\author{\vspace{-1em}}
\preauthor{}
\postauthor{}
\setlength{\droptitle}{-2cm}

\begin{document}
	\maketitle
	\section{Rozhodnětě a \underline{zdůvodněte}, jestli platí následující tvrzení:}
	\begin{enumerate}[A)]
		\item Pro funkci $f(n) = 4n^3 + 7n^2 + 2n + 11$ platí:
		\begin{align}
			f(n) \in \Omega(n), \\
			f(n) \in \Theta(n^3).
		\end{align}
		\vspace{2cm}
		\item Pro funkci $g(n) = n\log(10n) + n^2 + \sqrt{n + 1}$ platí:
		\begin{align}
			g(n) \in O(n^2), \\
			g(n) \in \Theta(n\log n).
		\end{align}
		\vspace{2cm}
		\item Pro funkci $h(n) = 3n^3 + 4^4n + 10^{10} $ platí:
		\begin{align}
			h(n) \in O(n^3) ,\\
			h(n) \in \Theta(n^3).
		\end{align}
	\end{enumerate}
	\vspace{2cm}
	Za dostatečné zdůvodnění se považuje výpočet odpovídající limity nebo ukázka toho, že lze nalézt $c_1, c_2$ a $n_0$ v definicích $\Omega, \Theta, O$, viz přednáška a doporučená literatura.
	Dostatečné zdůvodnění by mělo obsahovat komentář v přirozeném jazyce, alespoň na úrovni: ,,Tvrzení platí/neplatí protože:\dots``

	\section{Určete, která z následujících funkcí roste rychleji:}
	Doporučeno je postupovat výpočtem limit, případné invenci a alternativním postupům se ovšem meze nekladou, dokud jsou správné.
	Opět je očekáváno zdůvodnění typu: ,,Funkce $f$ roste rychleji, protože:\dots``
	\begin{enumerate}[A), itemsep=2.5em]
		\item $n^3$ nebo $n^4$
		\item $22n$ nebo $\log_2(n)$
		\item $n^{4}$ nebo $4^n$
		\item $n^4$ nebo $4n^3 + 4^{4^4}n^2 + 4^{4^{4^4}}$
		\item $n!$ nebo $n^n$
	\end{enumerate}
\end{document}
